\section{The complete phase-Fourier theorem}\label{sec:spectrum}

For $X,Y\in\Hc$, define
\begin{equation}\label{eq:def-energy-dft}
 \Ephase j=\norm{X+\ii^jY}^2,
 \qquad
 \Fmode k=\frac14\sum_{j=0}^3\ii^{kj}\Ephase j,
 \qquad 0\leq k\leq3.
\end{equation}

\begin{theorem}[Complete $C_4$ spectrum]\label{thm:spectrum}
With the conjugate-linear-first convention,
\begin{equation}\label{eq:complete-spectrum}
 \Fmode0=\norm X^2+\norm Y^2,
 \qquad \Fmode1=\ip YX,
 \qquad \Fmode2=0,
 \qquad \Fmode3=\ip XY.
\end{equation}
\end{theorem}

\begin{proof}
Put $S=\norm X^2+\norm Y^2$ and $c=\ip YX$.  Conjugate-linearity in the
first variable gives
\begin{align*}
 \Ephase j
 &=\ip{X+\ii^jY}{X+\ii^jY}\\
 &=S+\ii^j\ip XY+\ii^{-j}\ip YX
  =S+\ii^j\overline c+\ii^{-j}c.
\end{align*}
For every integer $m$, fourth-root orthogonality says
\begin{equation}\label{eq:root-filter}
 \frac14\sum_{j=0}^3\ii^{mj}
 =\begin{cases}1,&4\mid m,\\0,&4\nmid m.\end{cases}
\end{equation}
Substitution into~\eqref{eq:def-energy-dft} yields
\[
 \Fmode k=S\mathbf1_{k=0}+c\mathbf1_{k=1}
              +\overline c\mathbf1_{k=3}
 \quad (0\leq k\leq3),
\]
which is~\eqref{eq:complete-spectrum}.
\end{proof}

The theorem contains the familiar polarization identity as its $k=1$
coordinate, but it also displays the trivial mode, the identically zero
second mode, and the conjugate third mode.  Since every $E_j$ is real,
$F_3=\overline{F_1}$ is also the expected Fourier reality symmetry.

\begin{corollary}[Common-offset projection]\label{cor:offset}
Suppose a scalar $A$ is added independently of phase, so that
$E'_j=E_j+A$ for all four $j$.  Then
\[
 F'_0-F_0=A,
 \qquad
 F'_1-F_1=F'_2-F_2=F'_3-F_3=0.
\]
In particular, every phase-independent additive scalar contributes exactly
zero to the V59-selected mode $F_1$.
\end{corollary}

\begin{proof}
The increment is $\frac A4\sum_j\ii^{kj}$, and
equation~\eqref{eq:root-filter} applies.
\end{proof}

\begin{remark}[Type of the hypothesis]
Corollary~\ref{cor:offset} assumes the same scalar is proved to occur inside
each of the four labelled energies.  It does not license the replacement of
an unsigned estimate on a separate kernel by such a common term.  Establishing
that equality is part of the physical attachment problem.
\end{remark}
